\chapter{The Solow Model}\label{ch:b1:solow-model}

Countries that invest a large share of their output are, on the whole,
much richer than those that invest little. Among 181 countries, the quarter
that invested most over 1960--2019 --- at least $26.7\%$ of GDP on average --- had a
median income of 38,117 dollars a person in 2019 (in 2021 PPP dollars); the quarter that invested
least --- $15.4\%$ or less --- had 5,342. % figdata: solow_stats
Is investment the cause, and how much of a sevenfold gap can it explain?
This chapter takes the Solow model, met at school as a recurrence, into
continuous time, and confronts it with the data.

\begin{openingchart}{GDP per person in 2019 (2021 PPP dollars, log scale)
against the average investment share of GDP over 1960--2019 (the years
available, at least 30), 181 countries.}{Penn World
Table 11.0 (rgdpe, pop, csh\_i)}
\begin{tikzpicture}
\begin{axis}[omscatter, ymode=log, log ticks with fixed point, ytickten={3,4,5}, xmin=0, xmax=65, ymin=400, ymax=200000,
  xlabel={Average investment share, 1960--2019 (\% of GDP)}, ylabel={GDP per person, 2019}, width=0.7\linewidth,
  height=0.45\linewidth]
  \addplot[only marks, mark=*, mark size=1.3pt, omDef, fill opacity=0.6] table[x=s, y=y, col sep=comma] {figdata/bachelor-1/20-solow-model/out/invest_income.csv};
  \addplot[only marks, mark=none, nodes near coords, point meta=explicit symbolic, nodes near coords style={font=\scriptsize, anchor=west}]
    table[x=s, y=y, meta=country, col sep=comma] {figdata/bachelor-1/20-solow-model/out/invest_points.csv};
\end{axis}
\end{tikzpicture}
\end{openingchart}

The correlation between the investment share and the logarithm of income
is $0.5$. % figdata: solow_stats
Singapore invested $44.88\%$ of its output and reached 101,632 dollars a
person; Niger invested $16.03\%$ and had 1,447. % figdata: invest_points
The cloud is wide: China invested more than Norway, $33.62\%$ against
$28.96\%$, and had a quarter of its income. % figdata: invest_points

\begin{recall}
The Solow model as a recurrence, its \omterm{def:b1:solow-model:steady-state}{steady state} and convergence were met
at school; linear differential equations come from the mathematics
course. The \omterm{def:b1:long-run-model:long-run}{long-run model} of \cref{ch:b1:long-run-model} sets output
from given factors; this chapter lets capital grow.
\end{recall}

\section{The model in continuous time}\label{sec:b1:solow-model:continuous}

Output is $Y = F(K, L)$ with constant returns, so output per worker is
$y = f(k)$ with $k = K/L$ and $f(k) = F(k, 1)$. A constant share $s$ of output is
saved and invested, capital depreciates at the rate $\delta$, and the labour force
grows at the rate $n$:
\[
\dot K = sY - \delta K, \qquad \dot L = nL.
\]
Since $\dot k/k = \dot K/K - \dot L/L$,
\[
\dot k = s f(k) - (n + \delta)k.
\]
Investment per worker $sf(k)$ adds capital; \emph{break-even investment}
$(n + \delta)k$ is what is needed to equip new workers and replace worn
capital.

\begin{definition}[Inada conditions]\label{def:b1:solow-model:inada}
\index{Inada conditions}
A \omterm{def:b1:technology:production-function}{production function} per worker $f$ satisfies the \emph{Inada conditions}
when $f(0) = 0$, $f' > 0$, $f'' < 0$, $f'(k) \to \infty$ as $k \to 0$ and $f'(k) \to 0$ as $k \to \infty$.
\end{definition}

\begin{definition}[Steady state]\label{def:b1:solow-model:steady-state}
\index{steady state}
A \emph{steady state} is a capital per worker $k^* > 0$ at which $\dot k = 0$:
$sf(k^*) = (n + \delta)k^*$. Capital, output and consumption per worker are then
constant, and their totals grow at the rate $n$.
\end{definition}

\begin{proposition}[Existence and stability]\label{prop:b1:solow-model:stability}
If $f$ satisfies the \omterm{def:b1:solow-model:inada}{Inada conditions} and $0 < s < 1$, there is a unique \omterm{def:b1:solow-model:steady-state}{steady
state} $k^* > 0$, and from any $k(0) > 0$ capital per worker converges to it.
\end{proposition}

\begin{proof}
Let $\phi(k) = sf(k)/k$. It is decreasing (since $f$ is concave with $f(0) = 0$, $f(k)/k$
falls), tends to $\infty$ as $k \to 0$ (by l'Hôpital, $f(k)/k \to f'(0) = \infty$) and to $0$ as $k \to \infty$
(likewise $f(k)/k \to 0$). So $\phi(k) = n + \delta$ at exactly one $k^*$. Below $k^*$, $\phi(k) > n + \delta$
and $\dot k > 0$; above it $\dot k < 0$. Capital moves towards $k^*$ from either side,
and cannot overshoot it, since $\dot k = 0$ there.
\end{proof}

\begin{omfigure}
\begin{tikzpicture}
\begin{axis}[omrecord, grid=none, xlabel={Capital per worker $k$}, ylabel={Per worker},
  xmin=0, xmax=12, ymin=0, ymax=0.75, width=0.56\linewidth, height=0.4\linewidth, samples=80]
  \addplot[demand, domain=0:12] {0.2*x^(1/3)};
  \addplot[supply, domain=0:12] {0.06*x};
  \addplot[guide] coordinates {(6.086,0) (6.086,0.365)};
  \node[font=\footnotesize, anchor=south] at (axis cs:10.5,0.44) {$sf(k)$};
  \node[font=\footnotesize, anchor=south east] at (axis cs:9.3,0.6) {$(n+\delta)k$};
  \node[font=\footnotesize, anchor=south west] at (axis cs:6.15,0.07) {$k^* = 6.086$};
  \draw[->, thick] (axis cs:2,0.04) -- (axis cs:4.2,0.04);
  \draw[->, thick] (axis cs:11,0.04) -- (axis cs:8.6,0.04);
\end{axis}
\end{tikzpicture}

{\small The Solow diagram with $f(k) = k^{1/3}$, $s = 0.2$ and $n + \delta = 0.06$. Where
investment (blue) exceeds break-even investment (red), $k$ rises; where it
falls short, $k$ falls. The arrows point to the \omterm{def:b1:solow-model:steady-state}{steady state}.}
\end{omfigure}

\section{Comparative statics of the steady state}\label{sec:b1:solow-model:statics}

With $f(k) = k^{\alpha}$, the \omterm{def:b1:solow-model:steady-state}{steady state} solves $sk^{\alpha} = (n + \delta)k$:
\[
k^* = \left(\frac{s}{n + \delta}\right)^{1/(1 - \alpha)}, \qquad
y^* = \left(\frac{s}{n + \delta}\right)^{\alpha/(1 - \alpha)}.
\]
A higher saving rate raises the level of income per worker in the \omterm{def:b1:solow-model:steady-state}{steady
state}, but not its growth rate, which is zero; faster population growth
lowers the level. Taking logarithms,
$\partial \ln y^*/\partial \ln s = \alpha/(1 - \alpha)$. With a capital share of one-third the
elasticity is $\tfrac12$: doubling the saving rate raises steady-state income
by a factor $\sqrt 2$, about $41\%$.

That is far less than the opening chart seems to show. With $\alpha = 1/3$, the
saving rates of the top and bottom investment quartiles, $26.7\%$ and
$15.4\%$, predict an income ratio of $1.32$; the medians differ sevenfold.
% computed: code/bachelor-1/20-solow-model/exercises/solutions.py
Either other things differ with saving, or capital matters more than its
share of income suggests.

\section{The golden rule}\label{sec:b1:solow-model:golden}

Steady-state consumption per worker is $c^* = f(k^*) - (n + \delta)k^*$. A saving rate
of zero gives no capital and no consumption; a saving rate of one gives
capital and no consumption. In between there is a best.

\begin{definition}[Golden-rule capital stock]\label{def:b1:solow-model:golden-rule}
\index{golden rule}\index{golden-rule capital stock}
The \emph{golden-rule capital stock} $k_g$ is the steady-state capital per
worker that maximises steady-state consumption per worker.
\end{definition}

\begin{proposition}[The golden rule]\label{prop:b1:solow-model:golden}
$k_g$ satisfies $f'(k_g) = n + \delta$. With $f(k) = k^{\alpha}$, the saving rate that reaches it
is $s_g = \alpha$.
\end{proposition}

\begin{proof}
Every $k$ is a \omterm{def:b1:solow-model:steady-state}{steady state} for some $s$, namely $s = (n + \delta)k/f(k)$. So maximise
$c(k) = f(k) - (n + \delta)k$ over $k$: the first-order condition is $f'(k) = n + \delta$, and
$c'' = f'' < 0$. With Cobb--Douglas, $f'(k) = \alpha f(k)/k$, so $\alpha f(k_g) = (n + \delta)k_g = s_g f(k_g)$ and
$s_g = \alpha$.
\end{proof}

\begin{omfigure}
\begin{tikzpicture}
\begin{axis}[omrecord, grid=none, xlabel={Capital per worker $k$}, ylabel={Per worker},
  xmin=0, xmax=24, ymin=0, ymax=3.2, width=0.56\linewidth, height=0.4\linewidth, samples=80]
  \addplot[demand, domain=0:24] {x^(1/3)};
  \addplot[supply, domain=0:24] {0.06*x};
  \addplot[guide] coordinates {(13.095,0.786) (13.095,2.357)};
  \addplot[black, dashed, thin, domain=1:24] {2.357 + 0.06*(x - 13.095)};
  \node[font=\footnotesize, anchor=north west] at (axis cs:21,2.7) {$f(k)$};
  \node[font=\footnotesize, anchor=north west] at (axis cs:20,1.15) {$(n+\delta)k$};
  \node[font=\footnotesize, anchor=west] at (axis cs:13.3,1.6) {$c_g = 1.571$};
  \node[font=\footnotesize, anchor=north] at (axis cs:13.095,0.75) {$k_g$};
\end{axis}
\end{tikzpicture}

{\small The \omterm{def:b1:solow-model:golden-rule}{golden rule} with $f(k) = k^{1/3}$ and $n + \delta = 0.06$: consumption, the
vertical gap between output and break-even investment, is largest at
$k_g = 13.095$, where the slope of $f$ (the dashed tangent) equals $n + \delta$.}
\end{omfigure}

\begin{definition}[Dynamic inefficiency]\label{def:b1:solow-model:dynamic-inefficiency}
\index{dynamic inefficiency}
An economy is \emph{dynamically inefficient} when its capital per worker
exceeds the \omterm{def:b1:solow-model:golden-rule}{golden rule}, $k > k_g$: saving less would raise consumption at
every date.
\end{definition}

Below the \omterm{def:b1:solow-model:golden-rule}{golden rule} the choice is harder: saving more raises
consumption in the end but lowers it at first. Which path a society
should prefer depends on how it weighs the present against the future;
the Solow model, with $s$ given, cannot say.

\section{The speed of convergence}\label{sec:b1:solow-model:speed}

\begin{definition}[Speed of convergence, half-life]\label{def:b1:solow-model:speed}
\index{speed of convergence}\index{half-life}
Near the \omterm{def:b1:solow-model:steady-state}{steady state}, $\frac{d}{dt}\ln y \approx -\lambda(\ln y - \ln y^*)$. The rate $\lambda$ is the
\emph{speed of convergence}: the share of the remaining gap closed each
year. The \emph{half-life} $\ln 2/\lambda$ is the time it takes to close half of it.
\end{definition}

\begin{proposition}[Speed of convergence]\label{prop:b1:solow-model:lambda}
With $f(k) = k^{\alpha}$, $\lambda = (1 - \alpha)(n + \delta)$.
\end{proposition}

\begin{proof}
Write $\dot k/k = sk^{\alpha - 1} - (n + \delta) = s\,e^{-(1 - \alpha)\ln k} - (n + \delta)$. Linearise in $x = \ln k$ around
$x^* = \ln k^*$, where $sk^{*\alpha - 1} = n + \delta$: the derivative of the right side is
$-(1 - \alpha)sk^{*\alpha - 1} = -(1 - \alpha)(n + \delta)$. So $\dot x \approx -(1 - \alpha)(n + \delta)(x - x^*)$, and since
$\ln y = \alpha \ln k$, the same holds for $\ln y$.
\end{proof}

With $\alpha = 1/3$ and $n + \delta = 0.06$, $\lambda = 0.04$ and the \omterm{def:b1:solow-model:speed}{half-life} is $17.33$ years.
% computed: code/bachelor-1/20-solow-model/exercises/solutions.py
With Cobb--Douglas, the variable $z = k^{1 - \alpha}$ obeys an exact linear equation,
$\dot z = (1 - \alpha)[s - (n + \delta)z]$, so the gap in $z$ closes at the rate $\lambda$ exactly
(\cref{exo:b1:solow-model:exact}).

\begin{record}[Convergence at two per cent a year]
Across the states of the United States since 1880, the gap between the
typical poor and the typical rich state has narrowed at roughly $2\%$ a
year; across 73 regions of Western Europe since 1950, at about $2\%$ a year
too (Barro and Sala-i-Martin 1991). % ledger: b1:barro1991
The model with a capital share of one-third, $n + \delta = 0.06$ and \omterm{def:b1:technology:technical-progress}{technical
progress} of $0.02$ predicts $(2/3) \times 0.08 = 5.3\%$; it predicts $2\%$ only with a
capital share of $0.75$. % computed: code/bachelor-1/20-solow-model/exercises/solutions.py
Convergence happens, but more slowly than diminishing returns to physical
capital alone imply --- as if capital, broadly understood, had a larger
share.
\end{record}

\begin{omfigure}
\begin{tikzpicture}
\begin{axis}[omtime, xlabel={Years after the war}, ylabel={Output per worker}, xmin=0, xmax=80, ymin=1.4, ymax=1.9,
  width=0.56\linewidth, height=0.36\linewidth, samples=80]
  \addplot[demand, domain=0:80] {(3.3333 + (2.0999 - 3.3333)*exp(-0.04*x))^(0.5)};
  \addplot[guide, domain=0:80] {1.826};
  \addplot[guide] coordinates {(17.33,1.4) (17.33,1.648)};
  \node[font=\footnotesize, anchor=south east] at (axis cs:78,1.83) {$y^* = 1.826$};
  \node[font=\footnotesize, anchor=west] at (axis cs:18,1.5) {17.33 years};
\end{axis}
\end{tikzpicture}

{\small The weekend problem: Kestria's output per worker after half its
capital is destroyed, from $1.449$ back towards $1.826$, along the exact path.
After one \omterm{def:b1:solow-model:speed}{half-life} it is $1.648$.}
\end{omfigure}

\begin{record}[Postwar recoveries]
From 1950 to 1973 real GDP per person grew at $5.21\%$ a year in Germany and
$7.47\%$ in Japan, against $2.51\%$ in the United States and $2.4\%$ in the United
Kingdom. % figdata: postwar
In 1950 Germany's income per person was $0.478$ of the American and Japan's
$0.198$. % figdata: postwar
Economies far below their \omterm{def:b1:solow-model:steady-state}{steady state}, with their skills and
institutions intact and their capital destroyed, grew fastest --- as the
model predicts.
\end{record}

\section{Technical progress and balanced growth}\label{sec:b1:solow-model:progress}

Let technology be labour-augmenting, $Y = F(K, AL)$, with $A$ growing at the rate
$g$. Measure capital and output per \emph{effective worker}, $\tilde k = K/(AL)$
and $\tilde y = f(\tilde k)$:
\[
\dot{\tilde k} = sf(\tilde k) - (n + g + \delta)\tilde k.
\]
It is the same equation with $n + g + \delta$ in place of $n + \delta$.

\begin{definition}[Capital per effective worker, balanced growth path]\label{def:b1:solow-model:balanced}
\index{capital per effective worker}\index{balanced growth path}
\emph{Capital per effective worker} is $\tilde k = K/(AL)$. A \emph{balanced growth
path} is a path on which $Y$, $K$ and $C$ all grow at the same constant rate.
\end{definition}

\begin{proposition}[Balanced growth]\label{prop:b1:solow-model:balanced}
In the \omterm{def:b1:solow-model:steady-state}{steady state} of $\tilde k$, output, capital and consumption grow at the
rate $n + g$, output per worker at the rate $g$, and the capital--output ratio
$K/Y = s/(n + g + \delta)$ is constant.
\end{proposition}

\begin{proof}
$K = \tilde k^* AL$ and $Y = f(\tilde k^*)AL$ with $\tilde k^*$ constant, so both grow as $AL$, at
$n + g$; $C = (1 - s)Y$ too; $Y/L = f(\tilde k^*)A$ grows at $g$. From
$sf(\tilde k^*) = (n + g + \delta)\tilde k^*$, $K/Y = \tilde k^*/f(\tilde k^*) = s/(n + g + \delta)$.
\end{proof}

In the long run, then, income per worker grows only as fast as technology:
saving raises the level of the path, not its slope. Where $g$ comes from,
the model does not say.

\begin{definition}[Conditional convergence]\label{def:b1:solow-model:conditional-convergence}
\index{conditional convergence}\index{convergence, conditional}
\emph{Conditional convergence} is the prediction that economies grow faster
the further they are below their own \omterm{def:b1:solow-model:steady-state}{steady state}, which depends on their
saving, population growth and technology. It differs from \emph{absolute}
convergence, in which poor economies grow faster than rich ones whatever
their characteristics.
\end{definition}

\subsection*{The model to the data}

Taking logarithms of $y^*$ and adding human capital, the \omterm{def:b1:solow-model:steady-state}{steady state} gives
an equation that can be fitted across countries:
\[
\ln y_i = a + \frac{\alpha}{1 - \alpha}\,[\ln s_i - \ln(n_i + g + \delta)] + c \ln h_i,
\]
with $h_i$ a measure of human capital per worker. Fitting it by least squares
is the subject of \cref{ch:b1:simple-regression}; here only the results are
used.

\begin{measure}[Human capital in the Penn World Table]
The index \texttt{hc} turns a country's average years of schooling of the
population aged 15 and over into human capital per worker, $h = e^{\phi(S)}$,
where $\phi$ adds a return of $0.134$ for each of the first four years of
schooling, $0.101$ for years five to eight and $0.068$ for each further year,
following estimates of the wage gain from schooling. % ledger: b1:pwt-hc
It counts years in school, not what is learnt there.
\end{measure}

In the original fit on 98 countries, with averages over 1960--1985, the
equation with human capital explained $78\%$ of the variance of the
logarithm of income (adjusted $R^2$), with an implied capital share of $0.31$
and a human-capital share of $0.28$; the basic model without human capital
implied a capital share of about $0.6$ ($0.59$ in their intermediate sample;
Mankiw, Romer and Weil 1992).
% ledger: b1:mrw1992
Refitted on the Penn World Table for 145 countries, with averages over
1960--2019, income in 2019 and $g + \delta = 0.05$, the equation explains $66.6\%$ of the
variance, with an implied capital share of $0.445$; without human capital it
explains $46.3\%$, with an implied share of $0.645$. % figdata: solow_stats

\begin{omchart}{GDP per person in 2019 predicted by the augmented Solow model
against actual, 145 countries (2021 PPP dollars, log scales), with the 45-degree
line.}{Penn World Table 11.0 (rgdpe, pop, csh\_i, hc); fitted by the
tested script}
\begin{tikzpicture}
\begin{axis}[omscatter, xmode=log, ymode=log, log ticks with fixed point, xtickten={3,4,5}, ytickten={3,4,5},
  xmin=500, xmax=200000, ymin=500, ymax=200000, xlabel={Predicted}, ylabel={Actual}, width=0.56\linewidth,
  height=0.5\linewidth]
  \addplot[only marks, mark=*, mark size=1.3pt, omDef, fill opacity=0.6] table[x=fitted, y=y, col sep=comma] {figdata/bachelor-1/20-solow-model/out/mrw_fit.csv};
  \draw[gray, dashed] (axis cs:500,500) -- (axis cs:200000,200000);
  % figdata: mrw_points (fitted, actual): China (25890, 17770), India (8854, 7390),
  % Niger (1450, 1447), United States (49351, 68683)
  \begin{scope}[black!60, every node/.style={font=\scriptsize, black}]
    \draw (axis cs:49351,68683) -- (axis cs:14000,130000) node[anchor=east] {United States};
    \draw (axis cs:25890,17770) -- (axis cs:70000,9000) node[anchor=west] {China};
    \draw (axis cs:8854,7390) -- (axis cs:25000,2600) node[anchor=west] {India};
    \draw (axis cs:1450,1447) -- (axis cs:4500,750) node[anchor=west] {Niger};
  \end{scope}
\end{axis}
\end{tikzpicture}
\end{omchart}

The fit is good on average and poor for many countries. India's predicted
income is 8,854 dollars against 7,390 actual; the United States', 49,351
against 68,683. % figdata: mrw_points
Saving, population growth and schooling account for much of the ranking of
countries, but not for the size of the gaps at the extremes; what the
model leaves in the constant --- technology and efficiency --- is the subject
of \cref{ch:b1:growth-accounting}.

\begin{model}{The Solow model}
\assumes A closed economy; a constant saving rate; exogenous growth of
population $n$ and technology $g$; depreciation at the rate $\delta$; constant returns
and diminishing returns to capital.
\predicts Steady-state income per worker rising with $s$ and falling with $n$;
\omterm{def:b1:solow-model:conditional-convergence}{conditional convergence} to the \omterm{def:b1:solow-model:steady-state}{steady state} at the rate $(1 - \alpha)(n + g + \delta)$;
long-run growth of income per worker equal to $g$, whatever the saving rate.
\tested The augmented model explained $78\%$ of the cross-country variance
of log income with a capital share of $0.31$ (Mankiw, Romer and Weil 1992);
% ledger: b1:mrw1992
convergence across regions at about $2\%$ a year (Barro and Sala-i-Martin
1991). % ledger: b1:barro1991
\breaks Growth of technology $g$ is left unexplained; the basic model needs a
capital share near $0.6$ to fit; the measured \omterm{def:b1:solow-model:speed}{speed of convergence} is
slower than the model with a capital share of one-third predicts; and
capital does not flow from rich to poor countries as fast as diminishing
returns would imply.
\end{model}

\section*{Exercises}\label{sec:b1:solow-model:exercises}

\begin{exercise}[$\star$]\label{exo:b1:solow-model:steady}
With $f(k) = k^{0.3}$, $s = 0.25$, $n = 0.02$ and $\delta = 0.05$, find $k^*$, $y^*$ and $c^*$.
\end{exercise}

\begin{exercise}[$\star$]\label{exo:b1:solow-model:read}
From the opening chart's table, Singapore invested $44.88\%$ of GDP and had
an income of 101,632 dollars; China invested $33.62\%$ and had 17,770.
% figdata: invest_points
Does the comparison support the Solow model's prediction? What else
differs?
\end{exercise}

\begin{exercise}[$\star$]\label{exo:b1:solow-model:predict}
A country's saving rate rises permanently from $0.2$ to $0.25$. Predict the
growth rate of income per worker just after the change and in the long run.
\end{exercise}

\begin{exercise}[$\star$]\label{exo:b1:solow-model:golden-ex}
Show that with $f(k) = k^{\alpha}$ the golden-rule saving rate equals $\alpha$, and find
$k_g$ for $\alpha = 0.3$ and $n + \delta = 0.07$.
\end{exercise}

\begin{exercise}[$\star\star$]\label{exo:b1:solow-model:half}
With $\alpha = 1/3$ and $n + \delta = 0.07$, find the \omterm{def:b1:solow-model:speed}{speed of convergence} and the \omterm{def:b1:solow-model:speed}{half-life}.
\end{exercise}

\begin{exercise}[$\star\star$]\label{exo:b1:solow-model:ratio}
With $\alpha = 1/3$ and the same $n + \delta$, predict the ratio of steady-state incomes of
two countries saving $25\%$ and $10\%$. Compare with the sevenfold gap of the
investment quartiles in the opening chart.
\end{exercise}

\begin{exercise}[$\star\star$]\label{exo:b1:solow-model:balanced-ex}
With $n = 0.01$, $g = 0.02$, $\delta = 0.05$, $s = 0.2$ and $\alpha = 1/3$, give the growth rates of $Y$,
$Y/L$ and $K/Y$ on the \omterm{def:b1:solow-model:balanced}{balanced growth path}, and the value of $K/Y$.
\end{exercise}

\begin{exercise}[$\star\star$]\label{exo:b1:solow-model:inefficient}
An economy is in a \omterm{def:b1:solow-model:steady-state}{steady state} with $k^* > k_g$. Show that lowering the saving
rate to the \omterm{def:b1:solow-model:golden-rule}{golden rule} raises consumption per worker at every date.
\end{exercise}

\begin{exercise}[$\star\star$]\label{exo:b1:solow-model:fit}
In the refitted augmented model, India's predicted income in 2019 was
8,854 dollars and its actual income 7,390; the United States' were 49,351
and 68,683. % figdata: mrw_points
Compute each ratio of predicted to actual and say what the model leaves
out.
\end{exercise}

\begin{exercise}[$\star\star\star$]\label{exo:b1:solow-model:exact}
With $f(k) = k^{\alpha}$, set $z = k^{1 - \alpha}$. Show that $\dot z = (1 - \alpha)[s - (n + \delta)z]$, solve
it, and deduce that the \omterm{def:b1:solow-model:steady-state}{steady state} is globally stable.
\end{exercise}

\begin{exercise}[$\star\star\star$]\label{exo:b1:solow-model:linear}
For a general $f$ satisfying the \omterm{def:b1:solow-model:inada}{Inada conditions}, show by linearisation
that $\lambda = (1 - \alpha_k)(n + \delta)$, where $\alpha_k = k^*f'(k^*)/f(k^*)$ is capital's share at the
\omterm{def:b1:solow-model:steady-state}{steady state}.
\end{exercise}

\begin{exercise}[$\star\star\star$]\label{exo:b1:solow-model:refute}
What cross-country evidence would refute the Solow model's explanation of
income gaps between countries?
\end{exercise}

\section{Problem: Kestria after the War}\label{sec:b1:solow-model:problem}

\begin{problem}[{Weekend problem --- a war and a recovery}]\label{pb:b1:solow-model:kestria}
Kestria (invented) produces $y = k^{1/3}$ per worker, saves $s = 0.2$ of its
output, and has $n = 0.01$ and $\delta = 0.05$, with no \omterm{def:b1:technology:technical-progress}{technical progress} until
Part~IV. Time is in years.

\medskip\noindent\textbf{Part I --- The prewar steady state.}
\begin{enumerate}
  \item Find $k^*$.
  \item Find $y^*$ and $c^*$.
  \item Find the capital--output ratio.
  \item Find the \omterm{def:b1:technology:marginal-product}{marginal product} of capital, and compare it with $n + \delta$.
  \item Is Kestria above or below the \omterm{def:b1:solow-model:golden-rule}{golden rule}?
  \item Find $k_g$, $y_g$ and $c_g$.
\end{enumerate}

\medskip\noindent\textbf{Part II --- The capital stock halved.} A war destroys
half of the capital stock; the labour force is intact.
\begin{enumerate}[resume]
  \item Find $k$ and $y$ just after the war.
  \item Find the growth rate of $k$ just after the war.
  \item Find the growth rate of $y$ just after the war.
  \item Find the \omterm{def:b1:solow-model:speed}{speed of convergence} and the \omterm{def:b1:solow-model:speed}{half-life}.
  \item Write the exact path of $z = k^{2/3}$.
  \item Find $k$ and $y$ after one \omterm{def:b1:solow-model:speed}{half-life}.
  \item What share of the log gap $\ln y^* - \ln y$ has been closed after one \omterm{def:b1:solow-model:speed}{half-life}?
  \item What share of the gap in levels $y^* - y$?
  \item Why do the two shares differ from one half?
  \item Why do economies far below their \omterm{def:b1:solow-model:steady-state}{steady state} grow fast?
\end{enumerate}

\medskip\noindent\textbf{Part III --- A saving reform.} Back in the prewar \omterm{def:b1:solow-model:steady-state}{steady
state}, Kestria raises its saving rate to the \omterm{def:b1:solow-model:golden-rule}{golden rule}.
\begin{enumerate}[resume]
  \item What is the golden-rule saving rate?
  \item Find consumption per worker just after the reform.
  \item Find consumption per worker ten years after the reform.
  \item After how many years does consumption regain its prewar level?
  \item Find consumption in the new \omterm{def:b1:solow-model:steady-state}{steady state}.
  \item Who gains and who loses from the reform?
\end{enumerate}

\medskip\noindent\textbf{Part IV --- Technical progress.} Back at $s = 0.2$,
labour-augmenting technology now grows at $g = 0.02$.
\begin{enumerate}[resume]
  \item Find $\tilde k^*$ and $\tilde y^*$.
  \item Give the growth rates of output and of output per worker on the balanced path, and the capital--output ratio.
  \item Find the new \omterm{def:b1:solow-model:speed}{half-life}.
  \item State the result of Part II: without \omterm{def:b1:technology:technical-progress}{technical progress}, how long does Kestria take to recover half its gap after the war?
\end{enumerate}
\end{problem}
